\section{Discussion}
\label{sec:discussion}

\paragraph{Scope: synthetic tasks, $\leq$1M parameters.} Every experiment
in this paper is small by contemporary standards: the largest model tested
has $\approx183$K parameters ($d{=}64$), and every task is a synthetic
associative-memory or composition grammar with no natural-language content.
This project's own working practice treats results at this scale with
caution for good reason -- at a few hundred thousand parameters on natural
text, models barely learn unigram statistics, and no claim about
reasoning, generalization, or abstract thought is licensed by such a
model's behavior alone. That caution does not straightforwardly transfer
here, and it is worth being explicit about why. A behavioral claim on real
text (``this model reasons'') has no ground truth to check against except
more behavior, so a small model's poor performance is genuinely
uninformative about the mechanism in a larger model. The claims in
\S\ref{sec:taskd}--\ref{sec:taske} are different in kind: $\rank(Z)\geq K$
is a proven necessity for the exact solution at \emph{any} scale, so
finding effective rank $\approx K$, a causal step at $k\approx K$, and an
entity-subspace-restricted operator that is exactly the analytic target
are not behavioral proxies for a mechanism -- they are direct measurements
of a mechanism against a known, provable ground truth. What small scale
does still cost us is generality: we do not know whether the same
mechanism is recruited at 10M+ parameters, on real data, or under a
training objective with no comparably clean ground truth to check against.
That is precisely the transfer question \S\ref{sec:frontier}'s exactness
frontier and this project's own sequencing treat as the next, separate
step, gated on the results reported here -- not yet run.

\paragraph{What this does not show.} Task D and Task E are deliberately
abstract, chosen for provability rather than realism. They do not show
that matrix-valued rank helps performance on a real reasoning benchmark,
that the mechanism survives standard-scale pretraining, or that it competes
with any production architecture on accuracy or throughput. \S\ref{sec:related}'s
discussion of \citet{dziri2023faithfate} and \citet{wang2024grokked}
applies directly here: a mechanism shown to compose exactly on a synthetic
cycle with known structure is not evidence that the same mechanism
generalizes compositionally on tasks without that structure.

\paragraph{Concurrent, preliminary: DeltaNet.} A separate, still-running
campaign asks whether the same causal-rank story holds on a production
fast-weight architecture -- vanilla, single-Householder-per-token DeltaNet
\citep{schlag2021linear}, whose delta rule $S_t=S_{t-1}(I-\beta_tk_tk_t^\top)
+\beta_tv_tk_t^\top$ only ever writes within $\mathrm{span}(\{k_j\})$ by
construction. On the same Task-E-style composition grammar, the
unconstrained arm reaches $\mathrm{recovered\_frac@0.9}=1.00$ at every
tested hop (1--7, 21) at all four tested $(d,K)$ cells within
$2.5\times$ tier-default step budget, with zero training instability -- a
materially different outcome from the bespoke encoder used throughout this
paper, which plateaus sub-exact at the same $(d,K)=(64,32)$ operating
point after $150$K steps. Architecturally, entity-subspace rank equals
whole-matrix rank exactly, with no inflation: the delta rule has no
orthogonal complement for an optimizer to leave untouched, so
\S\ref{sec:taske-subspace}'s subspace-restriction fix is unnecessary here
by construction rather than by a fortunate optimization outcome. This is
consistent with, and architecturally cleaner than, this paper's Task E
result. It is \emph{not} yet a completed causal claim: the force-rank
staircase that would show a $k{=}K{-}1$ collapse (this paper's central test
in \S\ref{sec:taskd-m3} and \S\ref{sec:taske-depth}) is mid-campaign on
DeltaNet -- a $k{=}K$ (provably sufficient) cell is climbing but has not
reached ceiling within its tier budget, and the $k{=}K{-}1$ comparison has
not run. We report this paragraph only to bound what is already visible in
that campaign; the full DeltaNet result, including the causal staircase, is
a separate paper.

\paragraph{Falsification.} For the core causal claim
(\S\ref{sec:taskd-m3}, \S\ref{sec:taske-depth}): a force-rank-$(K{-}1)$
model reaching $\geq0.9\times$ ceiling accuracy at multiple $(d,K)$
operating points would falsify it. This has not happened -- the one
force-rank-$(K{-}1)$ Task E seed that converged at all showed exactly the
predicted depth-decay signature (Table~\ref{tab:taske-depth}), not ceiling
accuracy. For the exactness-frontier claim (\S\ref{sec:frontier}): a
$d\geq32$ cell reaching the pre-registered $>0.7$ bar at any tested budget
would revise it toward ``trainability, not exactness, is the limit''; this
has not happened through 100K--150K steps with confirmed-flat trailing
checkpoints. For the concurrent DeltaNet result: a completed Wave-B
force-rank grid showing no accuracy step between $k{=}K{-}1$ and $k\geq K$
even at extended budget would be the relevant negative there, distinct
from -- and, per \S\ref{sec:frontier}'s pattern, requiring the same
budget-extension discipline as -- every negative reported in this paper.

\paragraph{A standing caution about our own negative results.} Three
independent axes of this program each declared a cell dead at its first
tested budget before a $2$--$2.5\times$ extension reversed the verdict
(\S\ref{sec:frontier-pattern}). A reader evaluating any single ``wall''
claim in this literature, including ones in our own earlier project
documents that this paper's numbers supersede, should ask whether it was
re-tested at extended budget before being treated as final.
